\begin{hyperbolicbox}{$j$-invariant}
    \small
    \begin{defi}
        The $j$-invariant is a function from $\H$ to $\C$ defined by
        $$
        j(\tau)=1728\frac{g_2(\tau)^3}{\Delta(\tau)}
        $$
        The core property of the function is that $j(\tau)=j(\tau')$ if and only if there is a M\"obius map in $\operatorname{PSL}(2,\Z)$ such that $\tau'=\frac{a\tau+b}{c\tau+d}$, or equivalently, $\Lambda_\tau$ is homothetic to $\Lambda_{\tau'}$ \cite{silvermanAT}. Therefore, $j$-invariant acts as an isomorphism from $\H/\operatorname{PSL}(2,\Z)$ to $\C$. On the other hand, when we apply the correspondence between $\H/\operatorname{PSL}(2,\Z)$ and $\mathcal{L}/\sim$ we just deduced into the $j$-invariant, we can, in fact, view it as a parameterization of each equivalence class to points of $\C$.   
    \end{defi}
\end{hyperbolicbox}